\documentclass[11pt]{article}
\usepackage{amsmath,amssymb,amsthm,hyperref}
\newtheorem{proposition}{Proposition}
\newcommand{\E}{\mathbb E}
\newcommand{\Pp}{\mathcal P}
\newcommand{\B}{\mathcal B}
\title{A square-root local budget and code obstructions to conditional signed testing}
\author{Campaign 3395: Engineer continuation}
\date{November 4, 2026}
\begin{document}
\maketitle

\paragraph{Status and scope.}
The unrestricted fixed-factor lower bound is not proved.
The proposed matrix-general signed test is false, even with its
full-mean tolerance and exponential remainder retained.
After finding that obstruction, we also test a stronger statistic:
the coherent sparse-union defect of an actual balanced cut.
An explicit polynomial-size paired-group circuit defeats a
\emph{conditional-history} version of that stronger test.
Its fixed-source embedding is exponentially rare under the original
outer sampler, so it does not refute the actual-target signed estimate,
BP2, or the unrestricted lower bound.

The positive results are the square-root local budget and an explicit
conditional transfer to the target. Neither supplies the missing
global signed inequality. Reviewed V14 remains the strongest achieved
approximation theorem, at a factor tending to one.
No originality or publishable significance is certified.

We reuse reviewed K1--K2, CD5--CD7, J11, BP1--BP3 and BP8,
and the WP results without unchanged replay.
The code facts and all new inequalities below are proved directly.
In particular, no communication-protocol or external stability theorem
is imported. The elementary binary code construction is not claimed new.

\section{The local budget is valid}
For a nonnegative two-square matrix $Z$ put
\[
 \psi(Z)=
 \frac{(\sqrt{Z_{00}Z_{11}}-\sqrt{Z_{01}Z_{10}})^2}
      {\sum_{a,b}Z_{ab}},\qquad \psi(0)=0.                 \tag{ST1}
\]
Let $d_p(M)$ average this statistic at
$(u,v),(u,v\oplus b),(u\oplus a,v),(u\oplus a,v\oplus b)$,
where $u,v$ are independent uniform $k$-bit strings and $a,b$
are independent masks with Bernoulli-$p$ coordinates.

\begin{proposition}
If $R\ge Z$ entrywise is nonnegative and has rank at most one, then
\[
 \psi(Z)\le\sum_{a,b}(R_{ab}-Z_{ab}).
\]
For the actual cuts and CD5's unchanged law $\mu$, this implies
\[
 \sum_j\E_\mu\psi(Z_j)\le e^3b_Q,\qquad
 \sum_j\E_{\eta,z}\operatorname{mean}(M_j)
       =\E_\Delta Q\le3.                                 \tag{ST2}
\]
The means here are the \emph{full-cut} means, not the heavy means.
\end{proposition}
\begin{proof}
Set $A_{ab}=\sqrt{Z_{ab}}$, $B_{ab}=\sqrt{R_{ab}}$.
A nonnegative rank-one matrix has nonnegative outer factors;
taking their square roots shows that $B$ has rank at most one.
If $s_1\ge s_2$ are the singular values of $A$, then
\[
 \psi(Z)=\frac{s_1^2s_2^2}{s_1^2+s_2^2}\le s_2^2.
\]
Choose a unit vector $x$ in the kernel of $B$. Then
$\|A-B\|_{\rm F}^2\ge\|(A-B)x\|^2=\|Ax\|^2\ge s_2^2$.
Finally
$(\sqrt r-\sqrt z)^2\le r-z$ for $0\le z\le r$.
This proves the assertion, including the zero case.
The determinant identity used here is
$\det(AA^{\mathsf T})=(ad-bc)^2$ for
$A=\left(\begin{smallmatrix}a&b\\c&d\end{smallmatrix}\right)$.
The preserved bounded Jacobian calculation verifies this exact polynomial.

Use the actual K1 upper completion on each local union $U_j$.
If both row and column corners are distinct, K2 bounds the four
completion errors by $D_{Q_j}(U_j)$.
If either pair repeats, $\psi(Z_j)=0$; no repeated corner error
is charged twice. Sum over the actual cut and apply CD7.
For the second assertion, each conditional full-cut mean averages
back to $\E_\Delta Q_j$ under the same original law.
Sum $Q_j=Q$ before using $Q(X_\nu)\le3$.
There is no degree factor or gate-count multiplier on $b_Q$.
\end{proof}

\section{What a valid signed estimate would imply}
The proposed, subsequently refuted matrix assertion is
\[
 |\E_{u,v}(-1)^{u\cdot v}M(u,v)|
 \le\frac{\operatorname{mean}(M)}{64}
          +C d_p(M)+K\|M\|_\infty e^{-ck},               \tag{ST3a}
\]
with fixed positive constants, $k\ge256$,
$k\le q\le2^{111}k$, and $p=1/(2q)$.
Here and below the dot product in an exponent is over $\mathbb F_2$.

Had (ST3a) held, (ST2), $\|M_j\|_\infty\le3$,
$k_j\ge k_0=N/2^{134}$ on the original good event, and J11 would give
\[
 |\Lambda_\Delta(\Pp Q)|
 \le\frac3{64}+Ce^3b_Q+3Ks e^{-ck_0}
                  +3\sqrt{2s}e^{-N/2^{137}}.             \tag{ST3}
\]
The same conclusion follows from the actual-sampler alternative
(ST11) below; that alternative is not refuted here.
All conditional sign factors from unhonored bits have modulus one.
Also $q\le N/2^{23}\le2^{111}k_j$, so no sampler parameter is changed.

For an explicit rational upper bound $\overline C\ge C$, set
\[
 t=\frac1{6912(\overline C+1)},\qquad \kappa_0=1+t.
\]
Then $t<1/64$, $b_Q\le4t$, and $e^3<27$ imply
$Ce^3b_Q<1/64$. The non-remainder part of (ST3) is less than $1/16$.
For example, set
$a=\min(c/2^{135},2^{-137})>0$.
If $s<e^{aN}$, the remaining terms are at most
$3(K+2)e^{-aN/2}$, which is at most $3/16$ when
$N\ge(2/a)\log(16(K+2))$.
This contradicts BP3's lower bound greater than $7/16$.
Thus a proved (ST11) with explicit constants would give
$s\ge e^{aN}$ for all sufficiently large admissible $N$,
equivalently an exponential bound in $M_n=\Theta(N)$.
This is a conditional implication, not an achieved factor.

\paragraph{The earlier witnesses do not by themselves refute (ST3a).}
For rank-one $M$, $\psi$ vanishes and the Hadamard operator bound gives
$|\E(-1)^{u\cdot v}M|\le2^{-k/2}\sqrt{\E M^2}
\le2^{-k/2}\|M\|_\infty$.
For RG9's incidence matrix $M=J/2+B$ with density $1/t$,
regardless of its bit labeling,
$|\E(-1)^{u\cdot v}M|\le2^{-k-1}+1/t$.
For $t\ge256$, $2^k=t^2$, this is smaller than
$\operatorname{mean}(M)/64$, whose value is at least $1/128$.
Thus its former heavy-denominator obstruction is not transferred
to the deliberately different full mean.
In WP5, the unperturbed Boolean product is rank one across its cut,
and the perturbation has one-hot magnitude at most $\varepsilon$.
Its diverging exact-profile cost is not a diverging signed correlation.
Reviewed WP13 gives the family bound
$2b_Q+\exp(-N/2^{24})$ under the original sampler.
This observation is not a proof of (ST3a); it preserves both terms
that a support-only comparison would lose.

\section{A polynomial-density code defeats the matrix test}
Let $r\ge3$, $k=2^r-1$, and index the coordinates by the nonzero
vectors $h_i\in\mathbb F_2^r$. Define
\[
 C=\{u:\ \textstyle\sum_i u_i h_i=0\},\qquad
 \delta=|C|/2^k=1/(k+1).
\]
The parity-check map has rank $r$ because its columns include a basis.
No word of $C$ has weight one or two.
If $A_w$ counts its weight-$w$ words, deleting one marked coordinate
from a codeword determines that coordinate uniquely, as the sum of
the other columns. Hence
\[
 w A_w\le {k\choose w-1},\qquad
 A_w\le k^{w-1}/w!\quad(w\ge3).                          \tag{ST4}
\]

Put $M(u,v)=1_C(u)1_{\{u\cdot v=0\}}$ and $\mu=\operatorname{mean}(M)$.
Then
\[
 \|M\|_\infty=1,\quad
 \E(-1)^{u\cdot v}M=\mu
       =\frac1{2(k+1)}+2^{-k-1}.                         \tag{ST5}
\]
Indeed a nonzero row has exactly half its columns positive, and
the zero row has all columns positive.

For a Bernoulli-$p$ mask, write
$\rho_p=\Pr(a\in C\setminus\{0\})$.
Using (ST4), $kp\le1/2$, and the ratio at most $1/8$ between
successive terms of the exponential series after its cubic term,
\[
 \rho_p\le\frac1k\sum_{w\ge3}\frac{(kp)^w}{w!}
             \le\frac1{48k}\frac1{1-1/8}
             =\frac1{42k}.                              \tag{ST6}
\]
If one local row is outside $C$, it is zero.
If the two rows coincide, their minor is zero.
Thus a nonzero local statistic requires $u\in C$ and
$a\in C\setminus\{0\}$. Every zero-one two-square matrix has
$\psi\le1/2$: the only nonzero cases are two opposite ones
with value $1/2$, or three ones with value $1/3$.
Consequently
\[
 d_p(M)\le\tfrac12\delta\rho_p,\qquad
                    \frac{d_p(M)}{\mu}\le\frac1{42k}.    \tag{ST7}
\]
This bound holds for every $q\ge k$, not just the largest allowed $q$.

Fix arbitrary $C_*,K_*,c_*>0$ and take $q=k$.
Along $k=2^r-1$ with $r\ge9$, choose $k$ so large that
$C_* /(42k)\le1/64$ and
$2K_*(k+1)e^{-c_*k}\le1/64$.
The right side of (ST3a) is then at most $3\mu/64$,
whereas its left side is $\mu>0$.
This disproves the asserted existence of any fixed constants.
The density is polynomial, so the exponential remainder cannot
absorb this witness. This is an all-size proof, not a finite search miss.

\paragraph{Exact identities used in the replay.}
Every nonzero dual word has weight $(k+1)/2$. Character orthogonality,
obtained by summing $(-1)^{\ell\cdot\sum u_i h_i}$ over $\ell$,
therefore gives
\[
 \Pr(a\in C)=
  \frac{1+k(1-2p)^{(k+1)/2}}{k+1},\quad
 W_C(t):=\sum_w A_wt^w
 =\frac{(1+t)^k+k(1+t)^{(k-1)/2}(1-t)^{(k+1)/2}}{k+1}.
\]
Subtract $(1-p)^k$ for $\rho_p$.
For distinct nonzero rows $u,u'=u\oplus a$ in $C$,
averaging over $v$ and its mask gives the exact local value
\[
 \frac{7-3(1-2p)^{|u|}-3(1-2p)^{|u'|}
                    -(1-2p)^{|a|}}{48}.
\]
If one row is zero it is instead
$(1-(1-2p)^{|a|})/6$.
For the first formula the two row parities are independent uniform
bits. The mask parity patterns $10,01,11$ have conditional
mean statistics $1/6,1/6,1/4$, respectively; $00$ has zero.
The second formula uses the constantly positive zero row.

\section{Changing the test to coherent actual completions is not enough locally}
The next construction is stronger in its factor realization, but its
scope is still a conditional paired-group problem.
Use two two-choice row groups and two two-choice column groups
per coordinate. Write their variables as
$x_{i,0},x_{i,1}$, $x'_{i,0},x'_{i,1}$ and
$y_{i,0},y_{i,1}$, $y'_{i,0},y'_{i,1}$.
Set
\[
 r_u=\prod_i x_{i,u_i}x'_{i,u_i},\quad
 c_v=\prod_i y_{i,v_i},\quad d_v=\prod_i y_{i,v_i}y'_{i,v_i},
\]
\[
 P=\sum_{\substack{u,v\in C\\u\cdot v=0}}r_uc_v,\quad
 V=\sum_{v\in C}d_v,\quad
 A=\lambda F_S+P,\quad Q=AV,\quad \lambda=(k+1)^{-3}.    \tag{ST8}
\]
Here $S$ comprises all $2k$ row groups, exactly half the groups.
All selected labels can be designated nonroot labels.

\begin{proposition}
The polynomial $Q$ has a positive arithmetic DAG with at most
$64k(k+1)^2$ arithmetic gates. It can be ordered so that its sole
nonzero actual compulsory cut is $A$ with divided context exactly $V$.
It satisfies $Q\le(1+\lambda)F$ on the entire Boolean cube.
On the correlated bit corners its full cut matrix is
\[
 M(u,v)=\lambda1_C(v)+B(u,v),\qquad
 B(u,v)=1_C(u)1_C(v)1_{\{u\cdot v=0\}}.
\]
For $r\ge3$, with $\mu_B=\E B$,
\[
 \mu_B=\frac1{2(k+1)^2}+2^{-k-1},\qquad
 \E(-1)^{u\cdot v}M=\mu_B+\lambda2^{-k}
                      \ge\frac{32}{33}\E M.             \tag{ST9}
\]
For the actual sparse-union defect, not independently chosen local
completions,
\[
 \E_{\rm bits,masks}D_Q(U)
   =\mu_B(1+p)^{2k}
       \left[W_C\!\left(\frac{2p}{1+p}\right)-1\right]
   \le \frac{2\mu_B}{3k}\quad(q\ge k).                   \tag{ST10}
\]
\end{proposition}
\begin{proof}
Compute $P$ by a layered positive dynamic program whose state is
the two $r$-bit syndromes and the accumulated parity.
At coordinate $i$, the four choices $(u_i,v_i)$ multiply the
current value by $x_{i,u_i}x'_{i,u_i}y_{i,v_i}$ and update the
state by $(u_i h_i,v_i h_i,u_iv_i)$.
Accept the zero state. There are at most $2(k+1)^2$ states per layer.
Compute $V$ similarly with one syndrome and the paired column atoms.
Including their atoms, $F_S$, and the final three gates gives the
stated conservative gate bound. Sharing is counted in this original DAG.

Compute $V$ and $P$ before $A$. The small positive $\lambda F_S$
branch ensures $G(A)=S$; without it the optional column groups
would be compulsory and the asserted cut would be wrong.
$G(V)=\overline S$, and $A$ is the last balanced original gate.
The final product is full-group, and setting $A$ to zero kills
the whole output. The actual cut therefore has context $V$.

For Boolean inputs,
$\sum_u r_u\le F_S$ and
\[
 \left(\sum_v c_v\right)V
 \le\prod_i(y_{i,0}+y_{i,1})
               (y_{i,0}y'_{i,0}+y_{i,1}y'_{i,1})
 \le F_{\overline S}.
\]
Also $V\le F_{\overline S}$. Thus $PV\le F$ and
$\lambda F_SV\le\lambda F$, including every zero face.
The projection is exactly
$\lambda F_SV+\sum_{u,v\in C,\ u\cdot v=0}r_ud_v$.

For $r\ge3$ the dual code lies in $C$:
two distinct nonzero linear functionals on $\mathbb F_2^r$
are simultaneously one on $2^{r-2}$ columns, and a single one
is one on $2^{r-1}$ columns. These numbers are even.
Character orthogonality on $C$ gives
$\sum_{u,v\in C}(-1)^{u\cdot v}=|C||C^\perp|=2^k$,
proving the formula for $\mu_B$.
The other signed contribution is $\lambda2^{-k}$.
Since $\E M=\mu_B+\lambda/(k+1)$ and
$\lambda/(k+1)\le2\mu_B/(k+1)^2\le\mu_B/32$,
(ST9) follows.

Fix base strings and masks. Let ${\cal R},{\cal T}$ be the bit
subcubes spanned coordinatewise by the two rows and columns.
Let $n$ count the pairs of $B$ in ${\cal R}\times{\cal T}$ and
$m=|C\cap{\cal T}|$. Directly from the displayed projection,
\[
                         D_Q(U)=n(m-1).
\]
The low branch has no defect. If $C\cap{\cal T}$ is nonempty,
$m=L(b):=|C\cap\mathbb F_2^{\operatorname{supp}b}|$;
otherwise $n=0$. Hence $D_Q(U)=n(L(b)-1)$ always.
For fixed masks, averaging their base strings gives
$\E n=2^{|a|+|b|}\mu_B$.
Furthermore
\[
 L(b)-1=\sum_{w\in C\setminus\{0\}}
                 1_{\{\operatorname{supp}w\subseteq\operatorname{supp}b\}},
\]
and
$\E 2^{|b|}1_{\{\operatorname{supp}w\subseteq\operatorname{supp}b\}}
=(2p)^{|w|}(1+p)^{k-|w|}$.
This proves the equality in (ST10).
Here $k\,2p/(1+p)\le1$ and (ST4) give
$W_C(2p/(1+p))-1\le(1/k)\sum_{w\ge3}1/w!\le2/(9k)$.
The last bound uses the successive ratio at most $1/4$.
Finally $(1+p)^{2k}\le e<3$, proving the inequality.
\end{proof}

Thus even replacing $d_p$ by the actual conditional raw defect
does not yield a dimension-free version of (ST3a) for all these
full-group $Q\le3F$. The mean is polynomially small,
the relative defect tends to zero, and the original size is polynomial
in $k$. A polynomial-in-size exponential remainder still tends to zero
relative to the mean. This statement does not retain the actual
$R_n$ lower approximation promise or average over the original
pre-bit law, and is not asserted to refute a theorem that does.

\section{The original outer law removes this particular obstruction}
Embed the selected variables into any fixed original history with
these gadgets, and multiply by the other full-group factors.
Choose a balanced $S$ containing the selected row sources and
excluding the selected column sources. Remaining source groups can
fill the balanced sides. The $V$ factor requires all $2k$ specified
base column sources to be nonroot, independently of its codeword.
The one-hot bound is still $Q\le1+\lambda$.
Reviewed BP8 therefore gives directly
\[
       \E_\Delta Q\le(1+\lambda)(4/L)^{2k}.
\]
The bound needs neither a pin-cover expansion nor a count of codewords.
For $k$ proportional to $N$ it is exponentially small.
The high correlation in (ST9) was at a designated conditional
history, not under this full law.

One can preserve the complete actual-target promise formally:
put $P_0=\Pp Q$ (including the other full-group factors), and
$Q_H=H-\varepsilon P_0+\varepsilon Q$ for $H=F$ or $H=R_n$.
For $0<\varepsilon\le1/(64(1+\lambda))$ the baseline has nonnegative
coefficients, because $P_0\le(1+\lambda)F$ coefficientwise.
Then $\Pp Q_H=H$, $H\le Q_H\le(1+\varepsilon b_Q)H$,
and $b_{Q_H}=\varepsilon b_Q$.
This is a positive formal baseline, not a subtraction gate.
No polynomial-size representation of the $R_n$ baseline is supplied.
The added cut is still suppressed by the fixed nonroot-source event.
Consequently this lift also supplies no counterexample to BP2.

\paragraph{The changed, still unresolved approach.}
Retain the actual factors and use the full original outer law,
not a per-matrix or per-history theorem. Write
\[
 \begin{split}
 {\cal S}(Q)&=\sum_j\E_{\eta,z}1_{{\cal G}_j}
             |\E_{u,v}(-1)^{u\cdot v}M_j(u,v)|,\\
 {\cal D}_{\rm loc}(Q)&=\sum_j\E_\mu D_{Q_j}(U_j)\le e^3b_Q.
 \end{split}
\]
A sufficient remaining assertion, specifically for the actual
BP1 approximators, is
\[
 {\cal S}(Q)\le
 \frac1{64}\sum_j\E_{\eta,z}1_{{\cal G}_j}\operatorname{mean}(M_j)
       +C{\cal D}_{\rm loc}(Q)+3Ks e^{-ck_0}.            \tag{ST11}
\]
The conditional obstruction explains why the outer averaging cannot
be dropped even after coherent completion information is restored.
The source-rarity argument just proved handles this fixed-source
family only. It does not bound arbitrary overlapping optional
supports or their coefficient-weighted actual contexts.
No argument here establishes (ST11), and none excludes it.
This is the remaining proof-layer gap, not an unavailable tool.

\paragraph{Exact replay and its limits.}
\path{signed_test_check.py} checks the preserved typed Gram output,
binary minor values, code weight counts, exact mask-return formulas,
and the local statistic by independent complete enumeration at $k=7$.
It also checks an actual positive shared DAG at $k=3$, its compulsory
cut, formal projection, full-cube maximum by a proved row-pattern
reduction, and every conditional base/mask outcome with the original
corner marginals. The $k=3$ diagnostic uses the same construction but
does not invoke the $r\ge3$ dual-containment formula.
It has 172 original arithmetic gates, formal degree 15, one nonzero
actual cut, mean $13/256$, signed correlation $25/512$, and
average raw defect $343/124416$ at $p=1/6$.
Its full-cube defect maximum is exactly $1/8$, not a vanishing
whole-output defect inferred from its small local expectation.

For completeness, the full-cube reduction also permits different
choice sets $Y_i,Y'_i$ in the two column copies.
Let $n_u(Y)$ count the codewords allowed by $Y$ and orthogonal to $u$,
and let $m(Y\cap Y')$ count all codewords allowed by both copies.
The defect on an arbitrary row box is
\[
 \sum_{u\in C} r_u(X)
 \left[n_u(Y)m(Y\cap Y')-n_u(Y\cap Y')\right].
\]
Every bracket is nonnegative, and $\sum_u r_u(X)\le F_S(X)$.
The normalized maximum is therefore achieved by a fixed code row
and some column box. The replay covers all $2\cdot4^6=8192$
such cases at $k=3$, including zero faces.
The value $1/8$ is attained at row zero with every first column
copy allowing both labels and every second copy fixed at zero.
This proves that the reduced enumeration certifies the entire cube,
not only its paired bit corners.

Closed-form rational calculations also cover $k=511$, within the
matrix conjecture's dimension range; no enormous matrix is enumerated.
The small circuits are not simulations of the admissible large-$N$
graph family. The all-size conclusions follow from the proofs above,
not from promoting those finite diagnostics to general theorems.

Only the matrix-general testing route is retired.
The replacement (ST11), BP2, and the fixed-factor campaign remain open,
and all new mathematical judgments require Host-invoked independent review.
\end{document}
